% A second Write-Before-Read example, with a longer stretch of cycles in which
% the outputs are held without a read: two forwarded writes, then the delayed
% copy of the second one, then the plain hold.
%
% All values below are taken from a GHDL simulation of test/prog.asm, sampled
% around the instruction "SUB @R1++, @R1" at address 0x13C4 (label
% L_SUB_AM2_080).  R1 holds AM2_SRC1 = 0x1294 on entry and is post-incremented
% to AM2_SRC2 = 0x1295.

\documentclass{standalone}
\usepackage{timing}
\begin{document}
% Same drawing convention as write_before_read.tex.  Note that the register
% file's own RAM still answers 0xCD2E for R1 throughout t=1..t=4: every value
% shown on src_val_o/dst_val_o in those cycles comes from the forwarding path,
% not from the RAM.  The read issued at t=4 is the first one to land, at t=5.
\begin{wave}{8}{6}
 \nextwave{rd\_en\_i}    \bit{1}{1} \bit{0}{3} \bit{1}{2}
 \nextwave{src\_reg\_i}  \known{1}{1}    \unknown[x]{3}  \known{F}{2}
 \nextwave{src\_val\_o}  \unknown[x]{1}  \known{1294}{1} \known{1295}{3} \unknown[x]{1}
 \nextwave{dst\_reg\_i}  \known{1}{1}    \unknown[x]{3}  \known{1}{1}    \known{D}{1}
 \nextwave{dst\_val\_o}  \unknown[x]{1}  \known{1294}{1} \known{1295}{4}
 \nextwave{wr\_en\_i}    \bit{0}{1} \bit{1}{2} \bit{0}{3}
 \nextwave{wr\_reg\_i}   \unknown[x]{1}  \known{1}{2}                    \unknown[x]{3}
 \nextwave{wr\_val\_i}   \unknown[x]{1}  \known{1294}{1} \known{1295}{1} \unknown[x]{3}
\end{wave}

\end{document}
